\chapter{Observational Tests}
\label{chap:tests}

\section{Overview}

\SCT makes specific predictions that can in principle distinguish it from
$\Lambda$CDM and other standard frameworks. This chapter catalogs them in rough
order of how testable they are with current and near-future observational
capabilities.

\section{Tier 1: Already-observed potential signatures}

These are existing observations that may already be evidence for \SCT but are
conventionally interpreted otherwise.

\subsection{Galaxy rotation asymmetry}

The JWST JADES survey shows a $2\!:\!1$ asymmetry in galaxy rotation directions,
peaking near the Galactic poles, at $3.39\sigma$~\citep{shamir2025}.
\paragraph{\SCT prediction.} The universe inherits a preferred axis from the
parent's rotation. Galaxy rotation should show a hemisphere-based asymmetry
aligned with the parent's spin axis.
\paragraph{Status.} A hint, not a confirmation. Larger samples from ongoing JWST
surveys will sharpen this in the next few years. If the asymmetry persists at
higher significance and the Doppler-bias alternative is ruled out, this is direct
evidence for a parent-imprinted axis.

\paragraph{De-confounding (this work).} Running the test on the largest extant
catalogue---the clean spiral sample of Galaxy Zoo~1~\citep{lintott2011},
$30{,}126$ galaxies---shows how hard the measurement is (\Cref{fig:galaxydipole}).
There is a highly significant \emph{overall} handedness offset (clockwise
fraction $0.480$, a $3.9\%$ deficit, binomial $p\sim10^{-11}$), but it is the
known human perception bias~\citep{land2008}, not a sky dipole: a label-shuffle
null gives the dipole $p=0.62$ (no dipole beyond the monopole), and the best-fit
axis lies $7^\circ$ from the Galactic pole but $35^\circ$ from the CMB axis. On
the one large catalogue available, then, the asymmetry is systematics-dominated
and what little axis exists tracks the Milky Way, not the CMB. This does not
refute the parent-spin prediction---Galaxy Zoo~1 is a northern SDSS footprint
with uncontrolled handedness bias---but it pins down what a real test requires:
an all-sky, mirror-validated catalogue in which the discriminant is whether the
spin dipole favours the CMB axis over the Galactic pole. That is the single
measurement the framework most needs.

\begin{figure}[t]
  \centering
  \includegraphics[width=\textwidth]{galaxy_dipole}
  \caption{De-confounding the galaxy-spin axis on Galaxy Zoo~1 (clean spirals).
    \emph{Left:} the SDSS northern-cap footprint (red\,$=$\,clockwise,
    blue\,$=$\,counter); the fitted spin dipole (star) sits on the Galactic
    pole, $35^\circ$ from the CMB axis. \emph{Right:} the clockwise fraction sits
    a few percent below $0.5$ roughly independent of Galactic latitude---a global
    offset (perception bias), not a latitude systematic and not a significant
    dipole (shuffle null $p=0.62$).}
  \label{fig:galaxydipole}
\end{figure}

\subsection{``Axis of evil'' in the CMB}

Low-multipole CMB anomalies show suspicious alignments between the quadrupole,
octupole, and the ecliptic, at $\sim 2\text{--}3\sigma$.
\paragraph{\SCT prediction.} Parent rotation imprints anisotropy on the bounce
membrane, which manifests as low-multipole CMB anisotropies aligned with the
parent's rotation axis.
\paragraph{Status.} A long-standing anomaly. Conventional explanations
(foreground contamination, statistical fluke) have not fully resolved it. It
could be an \SCT signature.

\subsection{Evolving dark energy}

DESI year-1 BAO results favor $w_0 w_a$ models with $w_0 > -1$ and $w_a < 0$ over
strict $\Lambda$ at $\sim 2\text{--}3\sigma$~\citep{desi2024}, suggesting that
dark-energy density evolves.
\paragraph{\SCT prediction.} Dark energy reflects parent-accretion history, which
is generically non-constant. Specifically, observed evolution toward weaker dark
energy at recent times might indicate the parent has crossed into the
Hawking-dominated regime.
\paragraph{Status.} A statistical tension, not a confirmation. Years 2--5 of DESI
data will sharpen this. If evolving dark energy is confirmed at $>5\sigma$,
$\Lambda$CDM is in trouble and \SCT-like models become favored.

\subsection{Hubble tension}

A persistent $\sim 5\sigma$ disagreement between early-universe
($H_0 \approx 67$) and local ($H_0 \approx 73$) measurements.
\paragraph{\SCT prediction.} Bounce dynamics modify the relationship between
early and late expansion. A specific calculation is needed, but the framework has
degrees of freedom to potentially resolve the tension.
\paragraph{Status.} Calculational work needed. \SCT must reproduce both CMB
physics and local measurements consistently. If it can, this is a major
validation.

\subsection{$\sigma_8$ tension}

Lower-than-expected matter clustering inferred from low-redshift measurements
vs.\ CMB-extrapolated predictions.
\paragraph{\SCT prediction.} Dark matter is geometric (the parent's influence
through the membrane), not particulate. Its effective clustering behavior may
differ from CDM expectations.
\paragraph{Status.} Active observational research. \SCT-based dark-matter
modeling needed.

\subsection{JWST high-redshift massive galaxies}

JWST observations of unexpectedly massive galaxies at $z > 10$, challenging
standard structure-formation timelines.
\paragraph{\SCT prediction.} Bounce dynamics might allow earlier structure
formation than standard inflation, since the post-bounce expansion phase is
different.
\paragraph{Status.} Active research. Calculations of structure formation in
bounce cosmology needed.

\section{Tier 2: Near-future testable}

These can be tested with planned observational programs in the next 5--15 years.

\subsection{CMB spectral details}

If the CMB is parent Hawking radiation, the spectrum should be perfect Planck
(which it is, to $10^{-5}$). But subtle deviations or correlations might exist:
subtle spectral features from membrane geometry; correlations between CMB
temperature and dark-matter distribution; specific polarization patterns from a
rotating parent.
\paragraph{Tests.} Next-generation CMB experiments (CMB-S4, LiteBIRD) will
improve constraints on the CMB spectrum and polarization by orders of magnitude.

\subsection{Cosmic neutrino background structure}

The cosmic neutrino background (C$\nu$B) is predicted by standard cosmology but
not yet directly observed. If observed (PTOLEMY-class experiments): a pure-flip
framework predicts left-handed neutrinos in our convention; a strong-filter
framework predicts a fewer-than-expected antineutrino-like component; standard
cosmology predicts approximately equal $\nu$ and $\bar\nu$.
\paragraph{Tests.} Direct C$\nu$B detection is decades away with current
technology, but is being actively pursued.

\subsection{Primordial gravitational waves}

Inflation predicts a specific tensor-to-scalar ratio $r$. Bounce cosmology
predicts a different (typically smaller) $r$.
\paragraph{Tests.} LiteBIRD will constrain $r$ to the $\sim 10^{-3}$ level. If
observed at higher values, inflation is favored; if observed lower or absent,
bounce cosmology is favored.

\subsection{Dark-matter direct detection (negative result)}

\SCT predicts no fundamental dark-matter particles. WIMP, axion, sterile
neutrino, and other direct-detection experiments should continue to find
nothing.
\paragraph{Tests.} Ongoing programs (LZ, XENONnT, future generations) will probe
progressively lower cross-sections. Continued null results are weak evidence for
geometric (non-particulate) dark matter.

\subsection{Specific BAO and large-scale structure tests}

The relationship between baryon acoustic oscillations and dark-matter clustering
should follow specific \SCT-predicted patterns if dark matter is geometric.
Discrepancies with $\Lambda$CDM at large scales might be observable.
\paragraph{Tests.} DESI, Euclid, and the Roman Space Telescope will map
large-scale structure with unprecedented precision.

\section{Tier 3: Future testable, requires major advances}

\subsection{Gravitational decoherence experiments}

Tests of Di\'osi--Penrose gravitational collapse (or our preferred
reinterpretation, gravitational decoherence) at mesoscopic
scales~\citep{diosi1989,penrose1996}.
\paragraph{Tests.} Bouwmeester- and Aspelmeyer-class proposals using mesoscopic
mirror superpositions or matter-wave interferometers, within reach in the next
one to two decades. If gravitational decoherence is confirmed at predicted rates,
this strengthens the bridge between bounce cosmology and quantum mechanics that
we have argued for.

\subsection{Information-theoretic tests of dark matter}

If dark matter is entangled with the parent through the bounce, sufficiently
sensitive measurements might detect specific entanglement structures
(correlations between dark matter and Hawking-radiation-like backgrounds, etc.).
\paragraph{Tests.} Highly speculative. Requires precision quantum-cosmological
measurements not currently possible.

\subsection{Detection of parent-universe gravitational signatures}

If gravity propagates through the Cartasis membrane bidirectionally,
sufficiently long-wavelength gravitational waves might carry information about
the parent's mass distribution.
\paragraph{Tests.} Pulsar timing arrays (NANOGrav, etc.) probe nanohertz
gravitational waves. Detection of unexpected features at very low frequencies
might be parent signatures.

\section{Tier 4: Theoretical / computational tests}

These are not direct observations but consistency checks computable from the
framework.

\begin{itemize}
  \item \textbf{Filter efficiency yields the correct dark-matter ratio.} Compute
    filter efficiency from the Einstein--Cartan Lagrangian; predict the
    dark-matter-to-baryon ratio; compare to the observed $\sim 5\!:\!1$.
  \item \textbf{Bounce dynamics yield the correct CMB spectrum.} Compute the
    primordial perturbation spectrum from the torsion-driven bounce; compare to
    the observed CMB acoustic peaks.
  \item \textbf{Universe size distribution peaks near our parameters.} Simulate
    the supraverse population; show that observer-supporting universes cluster
    near our observed mass, expansion rate, and age.
  \item \textbf{Asymmetry magnitude from OG $+$ filter compounding.} Predict our
    observed baryon-to-photon ratio $\eta \sim 6 \times 10^{-10}$ from OG
    fluctuation statistics $+$ filter modulation across lineage depth.
  \item \textbf{Hubble tension resolution.} Show that bounce-modified
    early-universe dynamics produce $H_0$ consistent with both CMB and local
    measurements.
\end{itemize}

\section{What would falsify \SCT?}

\begin{enumerate}
  \item \textbf{No preferred axis in galaxy rotations at higher precision.} If
    Shamir's effect goes away with larger samples, the parent-rotation imprint is
    disfavored.
  \item \textbf{Dark-matter particle detection.} If a fundamental dark-matter
    particle is detected, the geometric interpretation fails.
  \item \textbf{Dark energy proven constant.} If DESI and successors confirm
    $w = -1$ at higher precision, dark-energy-as-parent-accretion fails.
  \item \textbf{CMB spectrum showing recombination-specific features inconsistent
    with a Hawking origin.}
  \item \textbf{Inflation-specific signatures confirmed.} If primordial
    gravitational waves are detected at inflation-predicted levels, the
    bounce-as-inflation-substitute is challenged.
  \item \textbf{CMB anomalies fully explained by foregrounds.} If the ``axis of
    evil'' and other anomalies are proven to be statistical or systematic
    artifacts, parent-rotation signatures are absent.
  \item \textbf{Bounce simulations fail to produce realistic cosmological
    evolution.} If detailed Einstein--Cartan bounce calculations cannot reproduce
    the CMB, structure formation, etc., the framework is wrong in detail even if
    the conceptual structure is appealing.
\end{enumerate}

\section{Priority ordering for tests}

For maximum scientific impact at minimum effort:

\begin{enumerate}
  \item \textbf{CMB spectrum from bounce} (Tier 4)---computable now, decisive if
    it works or fails.
  \item \textbf{Filter efficiency yields dark-matter ratio} (Tier 4)---computable
    now, strong validation if it works.
  \item \textbf{Shamir's effect with larger samples} (Tier 1)---observational
    program already running.
  \item \textbf{DESI dark-energy evolution} (Tier 1)---observational program
    already running.
  \item \textbf{Primordial gravitational waves from LiteBIRD} (Tier 2)---
    observational program in progress.
  \item \textbf{Hubble tension resolution} (Tier 4)---calculational.
  \item \textbf{Gravitational decoherence} (Tier 3)---experimental, near-future.
\end{enumerate}

Items 1, 2, and 6 are calculational work; items 3, 4, 5, 7 are observational
waits. The calculational work could be done by a small group in a few years; the
observational work happens on the experiments' timescales.
